% Chapter 2, typeset from lectures/L02: the README is the opener and the closing review, and the
% appendices are sections, with the same subsections: A, then the two toolchain walkthroughs C and
% D, then the review, and the exercises of Appendix B last.

\chapter[Hardware-Near Programming in C]{Hardware-Near Programming\\in C}
\label{c2:ch}

{\large\itshape No OS between you and the pins, only registers and the bits in them.\par}

\begin{chapterbox}{In this chapter}
\begin{itemize}
  \item Introduction: what's different about programming a microcontroller versus a normal
    computer.
  \item The board: Arduino Uno/Nano / ATmega328P, I/O ports and PIN$\leftrightarrow$PORT mapping.
  \item Bitwise operations: OR, AND, NOT, shifting, how register-level I/O actually works.
  \item Driving an output pin and reading an input pin (\code{DDRx}, \code{PORTx}, \code{PINx}).
  \item Two worked examples: LED follows a button; two LEDs blink while a button is held.
\end{itemize}
\end{chapterbox}

This chapter is where the book meets the hardware. It explains why embedded systems are programmed
close to the hardware, in C, without an OS underneath, and it teaches the one skill everything after
it is built on: reading and writing bitwise operations (\code{|}, \code{&}, \code{~}, \code{<<},
\code{>>}) to set, clear, and read individual bits in a register, so that you can configure a pin as
an input or an output and drive or read it correctly.

Read \secref{c2:sec:registers} first; it is the chapter's core material. The two sections after it
are for reference, and one of them is where you set up the toolchain the rest of the book uses:
\secref{c2:app:c} walks through flashing a Microchip Studio build onto the board via the Arduino
bootloader and \code{avrdude}, useful if you don't have a dedicated ISP programmer, and
\secref{c2:app:d} is the same idea from native Linux: installing \code{avr-gcc}/\code{avrdude},
compiling from the command line, and flashing over serial, including WSL2's USB-passthrough step.
Then work through the exercises in \secref{c2:sec:exercises}. The Reflection exercise is answered in
Appendix~\ref{app:answers}, and so are the values the expressions of Exercise~\ref{c2:ex:1} should
produce; the Code exercises have no printed solution.

% Section 2.1, from lectures/L02/appendix/a_registers_and_io.md.

\section{Registers and I/O}
\label{c2:sec:registers}\label{c2:app:a}

\subsection{What's different about embedded}
\label{c2:a:1}
\begin{itemize}
  \item A microcontroller (MCU) is a complete computer with no screen or keyboard; I/O happens
    through pins wired to sensors and actuators (buttons, LEDs, motors, sensors).
  \item Your code is still compiled to machine code, but there's no OS between you and the
    hardware; you write directly to memory-mapped registers.
  \item C is the default choice here: predictable, tiny memory footprint, close to the metal, and
    still far more productive than assembly. A very skilled assembly programmer can occasionally
    beat a modern compiler, but rarely by much, and it's almost never worth it.
\end{itemize}

\subsection{The board: Arduino Uno/Nano}
\label{c2:a:2}
Either board works for this book: same \textbf{ATmega328P} chip, same registers, same
PIN$\leftrightarrow$PORT mapping for D0--D13 and A0--A5. The Nano just adds two extra analog-only
pins (A6, A7) and drops the barrel power jack.

The ATmega328P is an 8-bit AVR microcontroller with three I/O ports, each covering a group of pins:

\begin{narrowtable}{@{}ll@{}}
  \thead{Board pin} & \thead{Register port} \\
  \midrule
  D0--D7 & \code{PORTD0}--\code{PORTD7} \\
  D8--D13 & \code{PORTB0}--\code{PORTB5} \\
  A0--A5 & \code{PORTC0}--\code{PORTC5} \\
\end{narrowtable}

\bookimage[0.3]{lectures/L02/appendix/images/arduino_uno_pinout.png}{The Arduino Uno: digital pins
D0--D13 along the top edge, analog pins A0--A5 along the bottom.}{c2:fig:uno}

\noindent Here's the Nano for direct comparison: different board outline and a Mini-USB connector
instead of the Uno's USB-B, but the same D0--D13/A0--A5 pin numbering and the same ATmega328P
underneath:

\bookimage[0.45]{lectures/L02/appendix/images/arduino_nano_pinout.png}{The Arduino Nano: the same
pin numbering, plus the analog-only pins A6 and A7.}{c2:fig:nano}

\noindent\code{PORTB} and \code{PORTD} are purely digital I/O. \code{PORTC} is the analog-capable
port (used with the ADC, Chapter~\ref{c4:ch}) but behaves as ordinary digital I/O when you're not
doing analog reads.

For consistency, every breadboard circuit diagram in this section (and in \secref{c2:app:c}) shows
an Arduino Uno. If you're using a Nano, wire it the same way: same chip, same D0--D13/A0--A5 pin
numbering, the only difference is the board outline and USB connector in the picture.

\subsection{Bitwise operations (this is the whole game)}
\label{c2:a:3}
There's no \code{digitalWrite()} here; you write directly to registers, and a register is just a
byte, so every pin operation is a bitwise operation on that byte.

\begin{booktable}{@{}p{5.4em}p{6.4em}L@{}}
  \thead{Operation} & \thead{C operator} & \thead{Effect} \\
  \midrule
  OR & \code{|} & sets one or more bits, leaves the rest (``bitwise addition'') \\
  AND & \code{&} & reads a bit, or (with NOT) clears one (``bitwise multiplication'') \\
  NOT & \code{~} & inverts every bit \\
  XOR & \code{^} & clears matching bits, sets differing bits \\
  NOR & \code{~(a | b)} & inverse of OR: bit set only where both operands are 0 \\
  NAND & \code{~(a & b)} & inverse of AND: bit set unless both operands are 1 \\
  XNOR & \code{~(a ^ b)} & inverse of XOR: bit set where operands match \\
\end{booktable}

\noindent There's no single C operator for NOR/NAND/XNOR; you get them by combining \code{~} with
\code{|}/\code{&}/\code{^} as shown above. Don't confuse NAND with the clear-bit idiom below:
\code{reg &= ~mask} is AND with a \emph{negated mask}, not a NAND of \code{reg} and \code{mask};
those are different operations that happen to both involve \code{~} and \code{&}:
\begin{itemize}
  \item Set bit 3 without touching the rest: \code{reg |= 0b00001000;}
  \item Clear bit 3 without touching the rest: \code{reg &= ~0b00001000;}
\end{itemize}

\subsubsection{Shifting}
\code{1U << n} puts a single \code{1} in bit position \code{n}; this is how you address ``which
pin'' without hand-writing the bit pattern:

\begin{ccode}
#define LED1 1U         // D9 -> PORTB1.
PORTB |= (1U << LED1);  // Enable LED1.
PORTB &= ~(1U << LED1); // Disable LED1.
\end{ccode}

\noindent\code{x << n} and \code{x >> n} are also multiply/divide by \code{2^n}, if you ever need
that framing.

\subsection{Driving an output, reading an input}
\label{c2:a:4}
\begin{itemize}
  \item \code{DDRx} (\textbf{D}ata \textbf{D}irection \textbf{R}egister) decides, per bit, whether a
    pin is an output (\code{1}) or input (\code{0}).
  \item On an output pin, \code{PORTx} sets the pin high/low.
  \item On an input pin, \code{PORTx} instead enables (\code{1}) or disables (\code{0}) the internal
    pull-up resistor. With it enabled, the pin reads a clean \code{1} when unconnected/floating and
    \code{0} when pulled to ground, instead of floating unpredictably.
  \item To read the actual input level, read \code{PINx} (not \code{PORTx}).
\end{itemize}

\begin{ccode}
#define BTN1 5U                               // D13 -> PORTB5
#define BTN1_PRESSED (PINB & (1U << BTN1)) // High if BTN1 is pressed, low otherwise.
\end{ccode}

\noindent The button circuit used in \secref{c2:a:5} and \secref{c2:a:6} doesn't actually rely on
this internal pull-up: it wires an \textbf{external pull-down} resistor instead (one leg of the
button to 5\,V, the other leg to D13 and, through a 10\,k\textohm{} resistor, to GND). That means
D13 reads \code{0} when idle and \code{1} when pressed, the opposite of the internal-pull-up pattern
described above. The code still enables the internal pull-up anyway (\code{PORTB = (1U << BTN1);}),
so you see \code{PORTx}'s input-mode behavior used directly, but it isn't what makes this circuit
work. It's harmless in practice: the internal pull-up and the external 10\,k\textohm{} pull-down
form a voltage divider, and with the internal resistor's typical value of around 35\,k\textohm{}
the idle pin sits near 1.1\,V, comfortably below the 1.5\,V ($0.3 \times 5$\,V) under which a pin
is guaranteed to read \code{0}. It isn't guaranteed, though: the datasheet only promises
20--50\,k\textohm, and at 20\,k\textohm{} the idle pin rises to about 1.7\,V, where the reading is
undefined. In a design of your own, don't enable a pull-up against an external pull-down.

\subsection{First program: LED follows a button}
\label{c2:a:5}
An LED is connected to D9, and a button is connected to D13. The button is continuously polled:
\begin{itemize}
  \item The LED is enabled if the button is pressed.
  \item The LED is disabled if the button isn't pressed.
\end{itemize}

\bookimage{lectures/L02/appendix/images/circuit_led_follows_button.png}{An LED on D9 and a button on
D13, with its 10\,k\textohm{} pull-down resistor to GND.}{c2:fig:ledbutton}

\begin{ccode}
#include <avr/io.h>

/** GPIO pins. */
#define LED1 1U // D9  -> PORTB1.
#define BTN1 5U // D13 -> PORTB5.

/** GPIO operations. */
#define LED1_ON PORTB |= (1U << LED1)         // Enable LED1.
#define LED1_OFF PORTB &= ~(1U << LED1)       // Disable LED1.
#define BTN1_PRESSED (PINB & (1U << BTN1)) // High if BTN1 is pressed, low otherwise.

/**
 * @brief Set up system.
 */
static void setup(void)
{
    // Configure LED1 as output.
    DDRB = (1U << LED1);

    // Configure BTN1 as input with its pull-up enabled.
    PORTB = (1U << BTN1);
}

/**
 * @brief Application entry point.
 *
 * @return 0 on termination of the program (should never occur).
 */
int main(void)
{
    setup();

    while (1)
    {
        // Read BTN1, enable LED1 if pressed, disable if not.
        if (BTN1_PRESSED) { LED1_ON; }
        else { LED1_OFF; }
    }
    return 0;
}
\end{ccode}

\noindent\code{LED1_ON}, \code{LED1_OFF}, and \code{BTN1_PRESSED} are the object-like macros covered
in \secref{c1:a:12}: plain text substitution, with no parameters. \code{BTN1_PRESSED} is used as a
value, inside a condition, so its whole body is parenthesized and it behaves correctly no matter
what surrounds it at the call site. \code{LED1_ON} and \code{LED1_OFF} are complete statements,
only ever used on their own, followed by a \code{;}.

You'll see a Doxygen comment above every non-trivial function throughout this codebase, written in
the imperative mood (``Configure\dots'', ``Compute\dots'', not ``Configures\dots'' or ``This
function configures\dots'').

\subsection{Blinking two LEDs, with a hand-rolled delay}
\label{c2:a:6}
LEDs on D9/D10 alternate every 200\,ms while the button on D13 is held; both off otherwise.

\bookimage{lectures/L02/appendix/images/circuit_two_led_blink.png}{Two LEDs on D9 and D10, and the
button on D13.}{c2:fig:twoleds}

\begin{ccode}
#define F_CPU 16000000UL // CPU frequency in Hz.

#include <stdint.h>

#include <avr/io.h>
#include <util/delay.h>

/** GPIO pins. */
#define LED1 1U // D9  -> PORTB1.
#define LED2 2U // D10 -> PORTB2.
#define BTN1 5U // D13 -> PORTB5.

/** GPIO operations. */
#define LED1_ON PORTB |= (1U << LED1)      // Enable LED1.
#define LED1_OFF PORTB &= ~(1U << LED1)    // Disable LED1.
#define LED2_ON PORTB |= (1U << LED2)      // Enable LED2.
#define LED2_OFF PORTB &= ~(1U << LED2)    // Disable LED2.
#define BTN1_PRESSED (PINB & (1U << BTN1)) // High if BTN1 is pressed, low otherwise.

/** Time parameters. */
#define TICK_1MS 1U         // Generate 1 ms tick in delay_ms().
#define BLINK_SPEED_MS 200U // Blink speed in ms.

/**
 * @brief Set up system.
 */
static void setup(void)
{
    // Configure LED1 and LED2 as outputs.
    DDRB = (1U << LED1) | (1U << LED2);

    // Configure BTN1 as input with its internal pull-up enabled.
    PORTB = (1U << BTN1);
}

/**
 * @brief Generate delay.
 *
 * @param[in] ms Delay duration in ms.
 */
static void delay_ms(const uint16_t ms)
{
    for (uint16_t i = 0U; i < ms; ++i)
    {
        _delay_ms(TICK_1MS);
    }
}

/**
 * @brief Blink the LEDs.
 *
 * @param[in] blink_speed_ms Blink speed in ms.
 */
static void led_blink(const uint16_t blink_speed_ms)
{
    // Enable LED1, disable LED2.
    LED1_ON;
    LED2_OFF;
    delay_ms(blink_speed_ms);

    // Disable LED1, enable LED2.
    LED1_OFF;
    LED2_ON;
    delay_ms(blink_speed_ms);
}

/**
 * @brief Turn both LEDs off.
 */
static void leds_off(void)
{
    LED1_OFF;
    LED2_OFF;
}

/**
 * @brief Application entry point.
 *
 * @return 0 on termination of the program (should never occur).
 */
int main(void)
{
    setup();

    while (1)
    {
        // Read BTN1, blink the LEDs if pressed, disable both LEDs if not.
        if (BTN1_PRESSED) { led_blink(BLINK_SPEED_MS); }
        else { leds_off(); }
    }
    return 0;
}
\end{ccode}

\noindent A few things worth noting:
\begin{itemize}
  \item \code{F_CPU} must be \code{#define}d \emph{before} including \code{<util/delay.h>}, since
    the delay functions compute their loop counts from the clock frequency at compile time.
    \code{int} is only 16 bits on this target, far too small for 16{,}000{,}000; the \code{UL}
    suffix makes the constant an \code{unsigned long} explicitly, rather than leaving its type to
    C's rules for unsuffixed constants.
  \item \code{<stdint.h>} is included explicitly for \code{uint16_t}, even though \code{<avr/io.h>}
    happens to pull it in already; include what you directly use rather than relying on a header's
    own includes, which can change.
  \item \code{_delay_ms()} only accepts a compile-time constant or macro as its argument, not a
    variable, so it can't take \code{ms} directly. \code{delay_ms()} works around this by calling
    \code{_delay_ms(TICK_1MS)} in a loop, incrementing a local counter each pass until it reaches
    \code{ms}.
  \item \code{_delay_ms()} (and by extension \code{delay_ms()}) \textbf{blocks} the entire program
    while it waits, fine for this example, but this exact limitation is what the timers chapter,
    Chapter~\ref{c5:ch}, exists to solve.
\end{itemize}

% Section 2.2, from lectures/L02/appendix/c_arduino_via_microchip_studio.md.

\section{Arduino via Microchip Studio}
\label{c2:sec:studio}\label{c2:app:c}

\subsection{Why this workflow}
\label{c2:c:1}
If you don't have a dedicated ISP programmer, you can still flash a Microchip Studio build onto the
board: the Arduino bootloader already sits on the chip, and the Arduino IDE already knows the exact
\code{avrdude} command line needed to talk to it over your board's COM port. The trick is simply to
borrow that command line and point it at Microchip Studio's \code{.hex} output instead of the
Arduino IDE's own. Everything below builds toward one reusable ``External Tool'' entry in Microchip
Studio that does exactly that.

\subsection{Wire the circuit}
\label{c2:c:2}
Same layout as the first example of \secref{c2:app:a}: LED anode on D9 (through a
220--470\,\textohm{} resistor to ground, 330\,\textohm{} is a good middle value), button on D13
(one leg to 5\,V, the other to ground through a 10\,k\textohm{} resistor, with D13 connected
between the button and the resistor).

\bookimage{lectures/L02/appendix/images/01_breadboard_wiring.png}{The circuit on the breadboard:
the LED on D9 and the button on D13, with its pull-down resistor.}{c2:fig:wiring}

\noindent Like the diagrams in \secref{c2:app:a}, this one shows an Arduino Uno. If you're using a
Nano instead: same chip, same pin numbering, just a different board outline in the picture.

\subsection{Build and test the sketch in Arduino IDE}
\label{c2:c:3}
This step uses the Arduino libraries (C++, an abstraction layer this book doesn't otherwise use)
purely to get a known-working upload before touching Microchip Studio. Open the Arduino IDE, start
a new sketch (\code{File → New}), and replace its contents with:

\begin{cppcode}
/** GPIO pins. */
#define LED1 9U  // D9.
#define BTN1 13U // D13.

/**
 * @brief Set up system.
 */
void setup()
{
    // Configure LED1 as output.
    pinMode(LED1, OUTPUT);

    // Configure BTN1 as input (the external 10 kOhm resistor pulls it down).
    pinMode(BTN1, INPUT);
}

void loop()
{
    // Read BTN1, enable LED1 if pressed, disable if not.
    if (digitalRead(BTN1)) { digitalWrite(LED1, HIGH); }
    else { digitalWrite(LED1, LOW); }
}
\end{cppcode}

\noindent This lights the LED on D9 while the button on D13 is held. It'll be rebuilt as real C in
\secref{c2:c:7}.

\code{pinMode()}, \code{digitalRead()}, and \code{digitalWrite()} are Arduino's abstraction over
the register-level I/O \secref{c2:app:a} covers directly: \code{pinMode(pin, mode)} sets a pin
\code{INPUT} or \code{OUTPUT} (a \code{DDRx} write), \code{digitalRead(pin)} returns
\code{HIGH}/\code{LOW} for a pin's current level (a \code{PINx} read), and
\code{digitalWrite(pin, value)} sets an output pin \code{HIGH} or \code{LOW} (a \code{PORTx} write).
Underneath, the Arduino core implements all three with exactly the register operations from
\secref{c2:app:a}, just hidden behind a function call and a pin-number lookup table.

Turn on verbose upload logging, since the next step depends on reading \code{avrdude}'s exact
command line out of it: \code{File → Preferences}, then check the box next to ``Show verbose output
during: Upload''.

\bookimage[0.95]{lectures/L02/appendix/images/02_file_menu_preferences.png}{The Arduino IDE's File
menu, with Preferences.}{c2:fig:c02}

\bookimage[0.95]{lectures/L02/appendix/images/03_preferences_verbose_upload.png}{The Preferences
dialog, with verbose output turned on for upload.}{c2:fig:c03}

\subsection{Install the board package, select the board and port}
\label{c2:c:4}
Install the AVR board package if it isn't already: \code{Tools → Board → Boards Manager}, then click
Install on ``Arduino AVR Boards by Arduino'' (skip this if it's already installed, as in the
screenshot below).

\bookimage[0.95]{lectures/L02/appendix/images/04_tools_board_boards_manager_menu.png}{Tools, Board,
Boards Manager.}{c2:fig:c04}

\bookimage[0.95]{lectures/L02/appendix/images/05_boards_manager_installed.png}{The Boards Manager,
with Arduino AVR Boards already installed.}{c2:fig:c05}

\noindent With the board connected, select it: \code{Tools → Board → Arduino AVR Boards → Arduino
Uno}. Then select its port: \code{Tools → Port}, and pick the COM port shown (e.g.\ COM4; the exact
number doesn't matter and can vary by machine).

\bookimage[0.95]{lectures/L02/appendix/images/06_tools_board_arduino_uno.png}{Selecting the Arduino
Uno as the board.}{c2:fig:c06}

\bookimage[0.95]{lectures/L02/appendix/images/07_tools_port_com4.png}{Selecting the board's port,
here COM4.}{c2:fig:c07}

\subsection{Upload and read back avrdude's command line}
\label{c2:c:5}
Click Upload in the top-left corner.

\bookimage[0.95]{lectures/L02/appendix/images/08_upload_button.png}{The Upload button, the arrow in
the top-left corner.}{c2:fig:c08}

\noindent In the Output window (the black pane), scroll up to the line reading ``Compilation
complete.'' and copy the \code{avrdude} command on the line beneath it; it should look something
like:

\begin{console}
"C:\Users\eripih2001\AppData\Local\Arduino15\packages\arduino\tools\avrdude\6.3.0-arduino17/bin/avrdude"
"-CC:\Users\eripih2001\AppData\Local\Arduino15\packages\arduino\tools\avrdude\6.3.0-arduino17/etc/avrdude.conf"
-v -V -patmega328p -carduino "-PCOM4" -b115200 -D
"-Uflash:w:C:\Users\ERIPIH~1\AppData\Local\Temp\arduino-sketch4CE86550DB8430968B53E52972DC68D2/Led1.ino.hex:i"
\end{console}

\bookimage[0.95]{lectures/L02/appendix/images/09_output_window_avrdude_command.png}{The avrdude
command in the Output window, on the line after ``Compilation complete.''}{c2:fig:c09}

\subsection{Turn it into a reusable command}
\label{c2:c:6}
This is the fiddliest part, but it's a one-time setup: reshaping that copied line into a
command/arguments pair Microchip Studio's ``External Tools'' feature can invoke directly.
\begin{enumerate}
  \item Open Notepad, paste the copied path in, and strip out all the quotation marks. Save it as
    \code{Arduino via Microchip Studio.txt}.
  \item After \code{bin/avrdude}, press Enter twice so the rest of the path lands two lines down.
  \item Add the extension \code{.exe} after \code{avrdude} on the top line.
\end{enumerate}

\bookimage[0.95]{lectures/L02/appendix/images/10_notepad_path_pasted.png}{The command in Notepad,
its quotation marks removed and the avrdude path on a line of its own.}{c2:fig:c10}

\bookimage[0.95]{lectures/L02/appendix/images/11_notepad_exe_appended.png}{\code{.exe} added after
\code{avrdude} on the top line.}{c2:fig:c11}

\begin{enumerate}[resume]
  \item Far to the right on the second line, just before \code{-Uflash}, press Enter so the rest
    lands on a new third line.
  \item On the two bottom lines, add quotes around the two file paths (the ones starting with
    \code{C:\\} and ending in \code{.conf} and \code{.hex} respectively). Afterward the two lines
    should look roughly like:
\end{enumerate}

\begin{console}
-C"C:\Users\eripih2001\AppData\Local\Arduino15\packages\arduino\tools\avrdude\6.3.0-arduino17/etc/avrdude.conf" -v -V -patmega328p -carduino -PCOM4 -b115200 -D
-Uflash:w:"C:\Users\ERIPIH~1\AppData\Local\Temp\arduino-sketch4CE86550DB8430968B53E52972DC68D2/Led1.ino.hex":i
\end{console}

\bookimage[0.95]{lectures/L02/appendix/images/12_notepad_arguments_split.png}{The arguments split
before \code{-Uflash}.}{c2:fig:c12}

\bookimage[0.95]{lectures/L02/appendix/images/13_notepad_quotes_added.png}{Quotes added around the
\code{.conf} and \code{.hex} paths.}{c2:fig:c13}

\begin{enumerate}[resume]
  \item Cut the file path on the bottom line (the part after \code{-Uflash}, between the quotes
    ending in \code{.hex}) and paste in \code{$(ProjectDir)Debug\$(TargetName).hex} instead:
    Microchip Studio expands these build-variable placeholders to the actual output path at run
    time, so this line no longer needs to change per project.
  \item Merge the two bottom lines back into one: click at the very start of the last line (right
    before \code{-Uflash}) and press Backspace. Only two lines should remain in total. Save the
    file.
\end{enumerate}

\bookimage[0.95]{lectures/L02/appendix/images/14_notepad_state_before_hex_replace.png}{The
\code{.hex} path replaced by Microchip Studio's build variables.}{c2:fig:c14}

\bookimage[0.95]{lectures/L02/appendix/images/15_notepad_hex_path_replaced.png}{The finished file:
the Command on the top line, the Arguments below it.}{c2:fig:c15}

\noindent You now have two lines: the \code{avrdude.exe} path (the \emph{Command}) and everything
else (the \emph{Arguments}): ready to paste into Microchip Studio's External Tools dialog in
\secref{c2:c:8}.

\subsection{Recreate the program in Microchip Studio}
\label{c2:c:7}
Open Microchip Studio and create a new \textbf{GCC C Executable Project}, named \code{Led1}.

\bookimage[0.95]{lectures/L02/appendix/images/16_new_project_dialog.png}{A new GCC C Executable
Project, named Led1.}{c2:fig:c16}

\noindent In the Device Selection dialog, choose \textbf{ATmega328P}.

\bookimage[0.95]{lectures/L02/appendix/images/17_device_selection.png}{Choosing the ATmega328P in
the Device Selection dialog.}{c2:fig:c17}

\noindent Replace the contents of \file{main.c} with the real-C equivalent of the sketch from
\secref{c2:c:3}, written to this codebase's conventions:

\begin{ccode}
#include <stdint.h>

#include <avr/io.h>

/** GPIO pins. */
#define LED1 1U // D9  -> PORTB1.
#define BTN1 5U // D13 -> PORTB5.

/** GPIO operations. */
#define LED1_ON PORTB |= (1U << LED1)         // Enable LED1.
#define LED1_OFF PORTB &= ~(1U << LED1)       // Disable LED1.
#define BTN1_IS_PRESSED (PINB & (1U << BTN1)) // High if BTN1 is pressed, low otherwise.

/**
 * @brief Set up system.
 */
static void setup(void)
{
    // Configure LED1 as output.
    DDRB = (1U << LED1);

    // Configure BTN1 as input with its pull-up enabled.
    PORTB = (1U << BTN1); // D13 as input, pull-up enabled
}

/**
 * @brief Application entry point.
 *
 * @return 0 on termination of the program (should never occur).
 */
int main(void)
{
    setup();

    while (1)
    {
        // Read BTN1, enable LED1 if pressed, disable if not.
        if (BTN1_IS_PRESSED) { LED1_ON; }
        else { LED1_OFF; }
    }
    return 0;
}
\end{ccode}

\noindent Then build it: \code{Build → Build Solution}.

\bookimage[0.95]{lectures/L02/appendix/images/18_mainc_build_menu.png}{Build, Build
Solution.}{c2:fig:c18}

\noindent The screenshots for this walkthrough were taken against earlier drafts of the program:
of the Arduino sketch in \secref{c2:c:3}--\ref{c2:c:5}, with other pin numbers, and of
\file{main.c}, with plain (non-\code{static}) functions and no Doxygen comments. The \emph{code
shown in the text} is what you should actually type in; the IDE mechanics (menus, dialogs, build steps) are otherwise
identical.

\subsection{Wire up the External Tool}
\label{c2:c:8}
Programming the chip externally requires the Advanced UI profile: \code{Tools → Select Profile},
choose \textbf{Advanced}, then \textbf{Apply}.

\bookimage[0.95]{lectures/L02/appendix/images/19_select_profile_menu.png}{Tools, Select
Profile.}{c2:fig:c19}

\bookimage[0.95]{lectures/L02/appendix/images/20_advanced_profile_dialog.png}{The Advanced
profile.}{c2:fig:c20}

\noindent Now create the External Tool entry: \code{Tools → External Tools}.

\bookimage[0.95]{lectures/L02/appendix/images/21_external_tools_menu.png}{Tools, External
Tools.}{c2:fig:c21}

\noindent In the dialog that appears, click \textbf{Add}. Title it \code{Arduino} followed by the
COM port used earlier in the Arduino IDE: e.g.\ \code{Arduino - COM4}. Paste the first line from
\code{Arduino via Microchip Studio.txt} into \textbf{Command}, and the second line into
\textbf{Arguments}. Check \textbf{Use Output Window}, then click \textbf{OK}.

\bookimage[0.95]{lectures/L02/appendix/images/22_external_tools_arduino_com4.png}{The External
Tools dialog with the Arduino - COM4 entry: its Command and its Arguments.}{c2:fig:c22}

\noindent Go back into \code{Tools} and select the tool you just created (e.g.\
\code{Arduino - COM4}).

\bookimage[0.95]{lectures/L02/appendix/images/23_tools_menu_arduino_com4_selected.png}{Running the
new tool from the Tools menu.}{c2:fig:c23}

\subsection{Flash and verify}
\label{c2:c:9}
If everything worked, the Output window should end with the number of bytes written, e.g.:

\begin{console}
Writing | ################################################## | 100% 0.05s

avrdude.exe: 194 bytes of flash written
avrdude.exe: safemode: lfuse reads as 0
avrdude.exe: safemode: hfuse reads as 0
avrdude.exe: safemode: efuse reads as 0
avrdude.exe: safemode: Fuses OK (E:00, H:00, L:00)

avrdude.exe done. Thank you.
\end{console}

\bookimage[0.95]{lectures/L02/appendix/images/24_output_window_flash_success.png}{A successful
flash, as the Output window reports it.}{c2:fig:c24}

\noindent The LED on D9 should now light up while the button on D13 is held, and stay off otherwise.
If the write fails instead, with something like:

\begin{console}
avrdude.exe: ser_open(): can't open device "\\.\COM4": Det går inte att hitta filen.
\end{console}

\noindent (the last part is Swedish Windows for ``The system cannot find the file specified''),
check that the board is connected, then check the Arduino IDE for its current COM port; if it's
changed, see below.

\subsubsection{If the COM port changes}
Reconnecting the board sometimes assigns a different COM port, which surfaces as the same
\code{ser_open()} error above. The fix is a second External Tool entry, identical to the first
except for the port: back in \code{Tools → External Tools}, add a new entry (e.g.\
\code{Arduino - COM3}), reusing the same Command, and an Arguments value with \code{PCOM4} swapped
for the new port (e.g.\ \code{PCOM3}).

\bookimage[0.95]{lectures/L02/appendix/images/25_external_tools_menu_for_com3.png}{Back to Tools,
External Tools, for a second entry.}{c2:fig:c25}

\bookimage[0.95]{lectures/L02/appendix/images/26_external_tools_arduino_com3.png}{The Arduino - COM3
entry: the same Command, with \code{-PCOM3} in its Arguments.}{c2:fig:c26}

\noindent Then select the new tool from \code{Tools} and run it the same way as before.

\bookimage[0.95]{lectures/L02/appendix/images/27_tools_menu_arduino_com3_selected.png}{Running the
COM3 tool from the Tools menu.}{c2:fig:c27}

% Section 2.3, from lectures/L02/appendix/d_building_and_flashing_on_linux.md. The recipe lines of
% the Makefile in the last subsection start with a TAB, as make requires.

\section{Building and Flashing on Linux}
\label{c2:sec:linux}\label{c2:app:d}

\subsection{Why this section}
\label{c2:d:1}
\Secref{c2:app:c} covers the Windows/Microchip Studio path. This is the native Linux equivalent:
the same \code{avr-gcc}/\code{avrdude} toolchain Microchip Studio and the Arduino IDE both wrap
under the hood, used directly from the command line. No IDE required.

This was written and tested from WSL2 (Ubuntu on Windows); the compiler/\code{avrdude} steps are
identical on any Debian/Ubuntu-based Linux, but the USB step of \secref{c2:d:5} is WSL2-specific;
skip it on a native Linux install, where the board just shows up as a device node once plugged in.

\subsection{Install the toolchain}
\label{c2:d:2}
Three packages: the AVR cross-compiler, its C standard library, and the flashing tool.

\begin{shell}
sudo apt -y update
sudo apt -y install gcc-avr avr-libc avrdude
\end{shell}

\noindent Verify:

\begin{shell}
avr-gcc --version
avrdude -v
\end{shell}

\subsection{Write the program}
\label{c2:d:3}
Same first program as \secref{c2:a:5}: LED on D9, button on D13, LED lit while the button is held.

\begin{ccode}
#include <avr/io.h>

/** GPIO pins. */
#define LED1 1U // D9  -> PORTB1.
#define BTN1 5U // D13 -> PORTB5.

/** GPIO operations. */
#define LED1_ON PORTB |= (1U << LED1)         // Enable LED1.
#define LED1_OFF PORTB &= ~(1U << LED1)       // Disable LED1.
#define BTN1_IS_PRESSED (PINB & (1U << BTN1)) // High if BTN1 is pressed, low otherwise.

/**
 * @brief Set up system.
 */
static void setup(void)
{
    // Configure LED1 as output.
    DDRB = (1U << LED1);

    // Configure BTN1 as input with its pull-up enabled.
    PORTB = (1U << BTN1);
}

/**
 * @brief Application entry point.
 *
 * @return 0 on termination of the program (should never occur).
 */
int main(void)
{
    setup();

    while (1)
    {
        // Read BTN1, enable LED1 if pressed, disable if not.
        if (BTN1_IS_PRESSED) { LED1_ON; }
        else { LED1_OFF; }
    }
    return 0;
}
\end{ccode}

\noindent Save it as \file{main.c}.

\subsection{Compile}
\label{c2:d:4}
Two steps: compile and link to an ELF, then extract the flash-able Intel HEX from it.

\begin{shell}
avr-gcc -mmcu=atmega328p -DF_CPU=16000000UL -Os -std=c11 -Wall -Wextra -o main.elf main.c
avr-objcopy -O ihex -R .eeprom main.elf main.hex
\end{shell}

\begin{itemize}
  \item \code{-mmcu=atmega328p} targets the right instruction set and register definitions.
  \item \code{-DF_CPU=16000000UL} matters the moment any code uses \code{<util/delay.h>}; harmless
    to always include.
  \item \code{-Os} (optimize for size) is the AVR default of choice; flash is measured in kilobytes.
  \item \code{avr-objcopy} strips the ELF down to the raw \code{.hex} \code{avrdude} actually writes
    to flash.
\end{itemize}
\code{avr-size main.elf} shows how much flash/RAM the build uses, worth checking once code grows
past a toy example.

\subsection{WSL2 only: attach the board's USB device}
\label{c2:d:5}
WSL2 doesn't share Windows' USB devices by default. \code{usbipd-win} bridges the two; this is a
one-time install, then a per-session attach.

\textbf{On Windows} (PowerShell, once): install the bridge:

\begin{shell}
winget install usbipd
\end{shell}

\noindent\textbf{On Windows} (PowerShell as Administrator, every time the board is (re)plugged in):

\begin{shell}
usbipd list                              # Find the board's BUSID.
usbipd bind --busid <BUSID>              # One-time per device.
usbipd attach --wsl --busid <BUSID>      # Do this each time you reconnect.
\end{shell}

\noindent\textbf{Back in WSL2}, confirm it showed up:

\begin{shell}
lsusb
ls /dev/ttyACM* /dev/ttyUSB* 2>/dev/null
\end{shell}

\noindent Arduino Uno-style boards (ATmega16U2 as USB-serial bridge) usually show up as
\file{/dev/ttyACM0}; boards using a CH340 or FTDI USB-serial chip (many Nano clones) show up as
\file{/dev/ttyUSB0} instead.

\subsection{Permissions}
\label{c2:d:6}
Without this, \code{avrdude} fails with a permissions error on the serial device. Add your user to
the \code{dialout} group (one-time), then start a new shell session for it to take effect:

\begin{shell}
sudo usermod -aG dialout $USER
\end{shell}

\noindent Log out and back in (or just open a new WSL terminal) for the group change to apply.

\subsection{Flash}
\label{c2:d:7}
Same idea as \secref{c2:app:c}: the board runs the Arduino bootloader, which \code{avrdude}'s
\code{arduino} programmer type talks to over serial.

\begin{shell}
avrdude -c arduino -p atmega328p -P /dev/ttyACM0 -b 115200 -U flash:w:main.hex:i
\end{shell}

\begin{itemize}
  \item \code{-c arduino} selects the STK500v1-over-serial protocol the Arduino bootloader speaks
    (this is the same programmer type the \code{avrdude} command of \secref{c2:app:c} uses via
    \code{-carduino}).
  \item \code{-p atmega328p} must match the \code{-mmcu} of \secref{c2:d:4}.
  \item \code{-P} is whatever device \secref{c2:d:5} (or, on native Linux, just plugging the board
    in) produced.
  \item \code{-b 115200} is the Optiboot-era default. Some older Nano boards ship a bootloader that
    only works at \code{-b 57600}; if 115200 fails with sync errors, that's the first thing to try
    instead.
\end{itemize}
A successful run ends with something like:

\begin{console}
avrdude: 194 bytes of flash written
avrdude: safemode: Fuses OK
avrdude done. Thank you.
\end{console}

\noindent The LED on D9 should now light up while the button on D13 is held, exactly like the
result of \secref{c2:app:c}, just reached without Windows or an IDE.

\subsection{Optional: a Makefile}
\label{c2:d:8}
Ties \secref{c2:d:4} and \secref{c2:d:7} together to replace the manual commands:
\begin{itemize}
  \item Call this file \file{Makefile}.
  \item Run \code{make} to compile the code.
  \item Run \code{make flash} to flash the firmware.
\end{itemize}
Adjust \code{PORT} for your setup:

\begin{makecode}
MCU   := atmega328p
F_CPU := 16000000UL
PORT  := /dev/ttyACM0
BAUD  := 115200

all: main.hex

main.elf: main.c
	avr-gcc -mmcu=$(MCU) -DF_CPU=$(F_CPU) -Os -std=c11 -Wall -Wextra -o $@ $<

main.hex: main.elf
	avr-objcopy -O ihex -R .eeprom $< $@

flash: main.hex
	avrdude -c arduino -p $(MCU) -P $(PORT) -b $(BAUD) -U flash:w:$<:i

clean:
	rm -f main.elf main.hex

.PHONY: all flash clean
\end{makecode}

% Section 2.4, from lectures/L02/README.md: the goals, the evaluation questions, and the next
% lecture. The README's Reference list, pointing at Appendices C and D, is in the chapter opener.

\needspace{10\baselineskip}
\section{Review}
\label{c2:sec:review}

\subsection*{What you should be able to do now}
After this chapter, you should:
\begin{itemize}
  \item Understand why embedded systems are programmed close to the hardware, in C, without an OS
    underneath.
  \item Be able to read and write bitwise operations (\code{|}, \code{&}, \code{~}, \code{<<},
    \code{>>}) to set, clear, and read individual bits in a register.
  \item Be able to configure a pin as input or output and drive/read it correctly.
\end{itemize}

\needspace{6\baselineskip}
\subsection*{Questions to test yourself}
You should be able to explain:
\begin{enumerate}
  \item What's the difference between \code{DDRx}, \code{PORTx}, and \code{PINx}, and when do you
    write to one versus read from another?
  \item How do you set bit 3 of a register without disturbing the other bits? How do you clear it?
  \item Why does an input pin need its pull-up resistor enabled, and what would you expect to happen
    if you left it disabled with nothing else connected to the pin?
\end{enumerate}

\needspace{6\baselineskip}
\subsection*{Looking ahead}
Chapter~\ref{c3:ch} is about interrupts:
\begin{itemize}
  \item Interrupts: reacting to events instead of polling for them.
  \item External interrupts (INT0/INT1) and Pin Change Interrupts.
\end{itemize}

% Section 2.5, from lectures/L02/appendix/b_exercises.md.

\section{Exercises}
\label{c2:sec:exercises}\label{c2:app:b}

\begin{aside}
\textbf{How to work these.} The Reflection exercise is answered on paper, from the chapter.
Exercise~\ref{c2:ex:1} needs no board: as its tip shows, a small program on your own computer checks
the expressions. The other Code exercises are programs for the board, and this is the chapter that
sets up the toolchain for them: build and flash them the way \secref{c2:app:c} (Microchip Studio)
or \secref{c2:app:d} (Linux) describes, or run them in the simulator of Appendix~\ref{sim:ch} if you
have no board yet.

\textbf{How to check your work.} The Reflection exercise, and the register values the expressions
of Exercise~\ref{c2:ex:1} should produce, are answered in Appendix~\ref{app:answers}. The Code
exercises have no printed solution: the board tells you whether a program does what the exercise
describes. Try each exercise before you read its answer.
\end{aside}

% ------------------------------------------------------------------------------------------
\exerciseset{Bit manipulation}
\label{c2:set:1}

\codeexercise{Set, clear and read a bit}
\label{c2:ex:1}
Given a register currently holding \code{0b10110010}, write single-line C expressions (using
\code{|}, \code{&}, \code{~}, \code{<<}) to:
\begin{parts}
  \item Set bit 0 without affecting any other bit.
  \item Clear bit 5 without affecting any other bit.
  \item Read whether bit 4 is set, as a boolean expression.
\end{parts}

\note[Tip] You don't need real hardware to test this; write a small program on your own machine and
check the results by invoking \code{printf()} from \code{<stdio.h>}:

\begin{ccode}
#include <stdint.h>
#include <stdio.h>

int main(void)
{
    const uint8_t reg = (1U << 3U);
    printf("%u\n", reg);
    return 0;
}
\end{ccode}

\codeexercise{LED macros}
\label{c2:ex:2}
A macro is defined as \code{#define LED 6U}. Write the two macros \code{LED_ON} and
\code{LED_OFF}, operating on a register called \code{PORTD}, following the OR/AND-NOT pattern used
in the chapter.

% ------------------------------------------------------------------------------------------
\exerciseset{Registers}
\label{c2:set:2}

\reflexercise{Which register?}
\label{c2:ex:3}
For each of the following, name the correct register (\code{DDRx}, \code{PORTx}, or \code{PINx})
and explain in one sentence why:
\begin{parts}
  \item Configuring D5 as an output.
  \item Reading whether a button on D2 is currently pressed.
  \item Enabling the internal pull-up resistor on D7 configured as an input.
\end{parts}

% ------------------------------------------------------------------------------------------
\exerciseset{Program}
\label{c2:set:3}

\codeexercise{Two buttons, two LEDs}
\label{c2:ex:4}
Wire up LEDs and buttons as described below:
\begin{itemize}
  \item \code{LED1} on D6 and \code{LED2} on D7.
  \item \code{BTN1} on D2 and \code{BTN2} on D3.
\end{itemize}
Write a complete program, from scratch, that:
\begin{itemize}
  \item Lights \code{LED1} if \code{BTN1} is pressed and turns it off otherwise.
  \item Lights \code{LED2} if \code{BTN2} is released and turns it off otherwise.
\end{itemize}
Don't copy the chapter's example directly; use these different pin numbers and re-derive the
macros yourself.

\codeexercise{Cycling three LEDs}
\label{c2:ex:5}
Wire up LEDs and buttons as described below:
\begin{itemize}
  \item \code{LED1} on D5, \code{LED2} on D6, and \code{LED3} on D7.
  \item \code{BTN1} on D4.
\end{itemize}
Write a complete program, from scratch, that:
\begin{itemize}
  \item While \code{BTN1} is held, cycles through the LEDs one at a time, in order (\code{LED1}
    $\rightarrow$ \code{LED2} $\rightarrow$ \code{LED3} $\rightarrow$ \code{LED1} \dots), lighting
    exactly one for 100\,ms before moving to the next.
  \item Turns off all LEDs otherwise.
\end{itemize}

